% Insert in main.tex in place of the existing numerical-method/results block.
% Keep this input within the existing non-TheoryOnly conditional block.
% Required main packages: amsmath, graphicx, booktabs, natbib, placeins, float.
% Add the supplied AvramFreddiGoreac2022 entry to references.bib.
\providecommand{\NumLatexRoot}{.}
\providecommand{\NumFigureRoot}{../figures/numerical_scenarios}
\providecommand{\NumTheoryRef}[2]{#2~\ref{#1}}
% MATLAB 解析公式评价；不含优化器结果。
\providecommand{\NumJRefEOne}{0.416186}
\providecommand{\NumJRefETwo}{1.529511}
\providecommand{\NumJRefEThree}{0.433981}
\providecommand{\NumJRefEFour}{1.268835}
\providecommand{\NumJRefEFive}{0.536233}
\providecommand{\NumSBRef}{0.54769512}
\providecommand{\NumSwitchSEOneRef}{0.51017169}
\providecommand{\NumSwitchIEOneRef}{0.13423054}
\providecommand{\NumMaxReleaseRef}{7.44746905}

% 从 data/numerical_scenarios 的真实结果生成；条件开关由 main.tex 定义。
\NumResultsVerifiedtrue
\providecommand{\NumHorizon}{300}
\providecommand{\NumIntervals}{4000}
\providecommand{\NumDegree}{2}
\providecommand{\NumMeshDetails}{容量重积分核查后，E2 使用 $N=8000$，E4 使用 $N=8000$；其余算例保留基准网格。图表均采用各例最终保存的结果。}
\providecommand{\NumJOclEOne}{0.416205}
\providecommand{\NumJOclETwo}{1.529523}
\providecommand{\NumJOclEThree}{0.433999}
\providecommand{\NumJOclEFour}{1.268831}
\providecommand{\NumJOclEFive}{0.536256}
\providecommand{\NumExitETwoStart}{5.85000000}
\providecommand{\NumExitETwoEnd}{5.88750000}
\providecommand{\NumExitEFourStart}{1.05000000}
\providecommand{\NumExitEFourEnd}{1.08750000}
\providecommand{\NumNeutralETwoCostDifference}{5.98410210273e-13}
\providecommand{\NumEFourSignedCostDifference}{-3.68275844664e-06}
\providecommand{\NumEFourCapacityExcess}{7.70286958773e-07}
\providecommand{\NumFindingWaiting}{实际控制输出中，E1 和 E2 分别识别出 $0\to 1\to 0$ 与 $0\to q_B\to 1\to 0$。按控制区间识别，E1 的首个完全跟踪区间从 $t\approx 7.050$ 开始，其状态约为 $(0.5072,0.1330)$，感染比例低于容量；E2 则出现正长度容量平台，随后转为完全跟踪。跳跃附近存在有限个过渡单元，因此这里的事件时刻只按当前网格精度解释。}
\providecommand{\NumFindingTrackingBoundary}{E3、E4、E5 的实际阶段依次为 $1\to 0$、$q_B\to 1\to 0$ 和 $1\to 0$。E3 与 E5 均从初始正长度区间开始完全跟踪；E4 先维持容量平台。E2、E4 的理论容量退出时刻均位于各自容量段与完全跟踪段之间的单网格过渡区间内，因此不将首个纯完全跟踪单元的起点直接等同于连续切换点。}
\providecommand{\NumCheckStatement}{对原始分段常数控制逐区间重积分后，五例最大容量超出量为 $\num{8e-07}$，与配点状态的最大差异为 $\num{8.6e-08}$；末尾 $20$ 个时间单位内最大控制幅值为 $\num{0}$。各例终端状态均通过零控制安全延拓核查。另对 E2 使用 $N=8000$ 加密，并对 E1 使用常值 $q=0.7$ 生成非理论初始猜测，所得解均通过同样核查；这些浮点检查用于排查截断及离散问题，不是连续可行性的严格证明。对 E2 另以常值 $q=0.7$ 及其积分状态为初猜重新求解，所得阶段顺序与解析初始化相同，成本绝对差为 $\num{5.98e-13}$，并通过同样的数值核查。}
\providecommand{\NumOverallFinding}{五例成本相对解析参考的最大差异约为 $0.00459\%$。计算结果支持这五个初值下的阶段顺序及切换位置预测；切换位置的偏差应结合控制区间宽度和过渡单元解释。E4 的离散成本比解析参考低约 $\num{3.68e-06}$，其重积分轨道同时存在约 $\num{7.7e-07}$ 的容量超出；该结果不能视为低于解析最优值的严格连续时间可行控制。}

\providecommand{\NumMeshDetails}{}

% Missing actual figures remain blank. Never substitute predicted OCP plots.
\newcommand{\NumFigureOrBlank}[2]{%
  \IfFileExists{\NumFigureRoot/#1}{%
    \includegraphics[width=0.97\linewidth]{\NumFigureRoot/#1}%
  }{%
    \fbox{\parbox[c][#2][c]{0.93\linewidth}{\mbox{}}}%
  }%
}
\newcommand{\NumFindingParagraph}[1]{\par\smallskip #1\par\smallskip}

\section{数值方法与实验设置}
\label{sec:numerical-method}

本节通过不同初值区域下的代表性算例，比较解析最优策略与直接配点求解结果，
重点考察等待阶段、容量弧和完全跟踪阶段的出现顺序及其切换位置。
算例组织借鉴 \citet[第6节]{AvramFreddiGoreac2022} 的区域相关情景设计，
但使用本模型的状态方程、控制成本及
\NumTheoryRef{def:regions}{分区定义}，不移用其控制结构。
全局最优性由\NumTheoryRef{thm:verification}{验证定理}给出；
数值计算用于检验解析结构的计算一致性。

所有主算例采用相同参数
\begin{equation}
 p=0.5,\qquad c=2,\qquad \gamma=0.3,\qquad K=0.15,
 \label{eq:numerical-common-parameters}
\end{equation}
仅改变初始状态。状态 $s,i$ 均为相对于初始总量的比例，允许 $s_0+i_0<1$。
相应地，$h=\gamma/(pc)=0.3$，$\ell=r=0.15$。
这些参数用于说明理论结构，不视为对某一具体疫情的参数标定。

数值优化保持原动力学和线性目标：
\begin{equation}
\begin{aligned}
 \min_{q}\quad &J_T(q)=\int_0^T pcq(t)\,\mathrm dt,\\
 \text{满足}\quad
 &\dot s=-c[p+(1-p)q]si,\qquad
 \dot i=[pc(1-q)s-\gamma]i,\\
 &q\in[0,1],\quad 0\le s\le1,\quad 0\le i\le K,\quad
 (s(0),i(0))=(s_0,i_0).
\end{aligned}
\label{eq:numerical-ocp}
\end{equation}
不添加二次控制正则化，也不预设阶段数、切换时刻或容量段控制。
因此，容量控制 $q_B(s)=1-h/s$ 及其离开位置是待与优化输出比较的理论预测，
而不是数值优化中的额外阶段约束。

理论控制问题定义在无限时间区间上，而直接配点计算采用一个足够长的固定时域
$[0,T]$ 进行截断近似。数值问题不增加终端状态约束或终端成本，
也不要求终点预先落入解析构造的某个区域。
这种处理与 \citet[第6节]{AvramFreddiGoreac2022} 的数值展示思路相同：
选取足够长的固定时域，使主要控制阶段在终止时刻之前已经结束，
再比较有限时域数值解与解析结构。这里不把“$T$ 足够大”解释为
有限时域问题与无限时域问题在一般意义下的严格等价性。

作为截断效应的轻量后处理，每个被报告的算例在求解后检查两点：
控制是否已在终点前的一段时间内回到 $q\approx0$，以及终端状态从此继续取
$q=0$ 时的未来无控制峰值是否不超过 $K$。该检查不进入优化约束，
不预设解除状态或解除时刻，也不单独设置图表。

\ifNumResultsVerified
计算采用
\else
计算拟采用
\fi
MATLAB/OpenOCL 直接配点方法\citep{Koenemann2019OpenOCL}，
通过 CasADi\citep{Andersson2019CasADi} 构造离散问题，基准网格将 $[0,T]$ 划分为
$N=\NumIntervals$ 个控制区间，设置 $T=\NumHorizon$、$d=\NumDegree$，
控制在各区间内为常数，离散非线性规划由 IPOPT 采用局部优化方法求解\citep{WachterBiegler2006}。
数值收敛用于复核解析结构，不承担连续问题全局最优性的证明。
\NumMeshDetails
解析轨道按正文的自然段、容量段和完全跟踪段公式独立生成；
用作初始猜测时仍不固定其阶段结构。
数值控制成本按完整求解时域计算为
\begin{equation}
 J_{\mathrm{OCL}}=pc\sum_{k=1}^{N}q_k\Delta t_k,
 \label{eq:numerical-cost-sum}
\end{equation}
不以绘图窗口作为积分范围。
优化输出的容量约束、长时域尾段及终端无控制延拓在结果采用前另行核查。
% Fill from actual checks only. Empty until the experiments have run.
\NumFindingParagraph{\NumCheckStatement}

\section{不同初值区域下的数值验证}
\label{sec:numerical-results}

\subsection{代表初值与相平面分区}

选取表\ref{tab:numerical-scenarios}中的五个初值，分别覆盖两种等待区、
内部立即完全跟踪区以及容量线上内生阈值两侧的状态。
算例 E1 与 E2 保持 $i_0$ 相同，用来区分两种等待方式；
算例 E4 与 E5 保持 $i_0=K$，用来检验容量初值是否需要先维持平台。
算例 E3 补充了 $i_0<K$ 时直接采用完全跟踪的情形。

由安全边界和切换条件的数值评价得
\begin{equation}
 s_B\approx\NumSBRef,
 \label{eq:numerical-sB}
\end{equation}
其中 $(s_B,K)$ 为切换曲线与容量线的交点。
图\ref{fig:numerical-regions}用同一相平面标记初值、区域和相应轨道。
两种等待分支之间的边界使用正文定义的自然轨道标号
$\Phi=\Phi_B$，不是整条竖线 $s=s_B$。
所有区域均限制在 $s>0$、$0<i\le K$、$s+i\le1$ 的物理域内。

\begin{table}[H]
\centering
\small
\setlength{\tabcolsep}{4.2pt}
\caption{代表初值、理论控制结构与控制成本。$J_{\mathrm{ref}}$ 为解析公式的数值评价，
$J_{\mathrm{OCL}}$ 为直接配点返回控制在完整求解时域内的成本。}
\label{tab:numerical-scenarios}
\begin{tabular}{@{}clclrr@{}}
\toprule
算例 & $(s_0,i_0)$ & 区域 & 理论结构 & $J_{\mathrm{ref}}$ & $J_{\mathrm{OCL}}$\\
\midrule
E1 & $(0.75,0.01)$ & $\mathcal W_\Gamma$ & $0\to1\to0$ & \NumJRefEOne & \NumJOclEOne\\
E2 & $(0.99,0.01)$ & $\mathcal W_K$ & $0\to q_B\to1\to0$ & \NumJRefETwo & \NumJOclETwo\\
E3 & $(0.50,0.14)$ & $\mathcal T$ & $1\to0$ & \NumJRefEThree & \NumJOclEThree\\
E4 & $(0.70,0.15)$ & $\mathcal B$ & $q_B\to1\to0$ & \NumJRefEFour & \NumJOclEFour\\
E5 & $(0.50,0.15)$ & $\mathcal T\cap\{i=K\}$ & $1\to0$ & \NumJRefEFive & \NumJOclEFive\\
\bottomrule
\end{tabular}
\end{table}

\begin{figure}[H]
\centering
\NumFigureOrBlank{FigN1_regions.pdf}{4.2cm}
\caption{不同初值区域及代表轨道的相平面对照。
图中标出安全边界、切换曲线 $\Gamma_K$、容量线 $i=K$、两种等待区的公共边界，
以及 E1--E5 的初始状态。理论几何与直接配点轨道采用不同线型或标记。}
\label{fig:numerical-regions}
\end{figure}

\subsection{两种等待区域的对照}

算例 E1 的初值 $(0.75,0.01)$ 属于 $\mathcal W_\Gamma$。
根据\NumTheoryRef{def:candidate}{反馈构造}，先取 $q=0$，
沿自然轨道到达 $\Gamma_K$ 后采用完全跟踪，并在进入安全集后解除。
理论结构为
\[
 0\longrightarrow1\longrightarrow0.
\]
解析公式给出的完全跟踪起点约为
\[
 (s_\Gamma,i_\Gamma)
 \approx(\NumSwitchSEOneRef,\NumSwitchIEOneRef),
 \qquad i_\Gamma<K.
\]
因此，本例的关键检验是数值优化是否在容量尚未达到时开始完全跟踪，
以及是否没有出现正长度容量平台。

算例 E2 的初值 $(0.99,0.01)$ 属于 $\mathcal W_K$。
理论预测先自然达到容量，再采用 $q_B(s)=1-h/s$，
沿 $i=K$ 运行至 $(s_B,K)$，之后转入完全跟踪并在安全边界解除。
其结构为
\[
 0\longrightarrow q_B(s)\longrightarrow1\longrightarrow0.
\]
图\ref{fig:numerical-waiting}并列比较这两种等待方式。
平台是否出现、完全跟踪开始时的感染比例及切换点的位置，
共同构成本组情景的比较内容。

\begin{figure}[H]
\centering
\NumFigureOrBlank{FigN2_waiting.pdf}{5.0cm}
\caption{两种等待区域下的控制与感染比例：左列为 E1，右列为 E2；
上排为 $q(t)$，下排为 $i(t)$。
每幅图比较解析参考与 OpenOCL 结果，数值控制按分段常数形式显示，
感染图中的水平线表示容量 $K$。}
\label{fig:numerical-waiting}
\end{figure}

% EXPERIMENT_REQUIRED: actual E1/E2 observations and a few useful event values.
\NumFindingParagraph{\NumFindingWaiting}

不存在容量平台与全时域内不接触容量并非同一命题。
解析策略在安全边界解除后，感染比例可以沿无控制轨道再次上升至 $K$，
但不超过它；此时 $q=0$，不应将这种后续接触误判为容量控制弧。

\subsection{立即完全跟踪与容量边界初值}

算例 E3 的初值 $(0.50,0.14)$ 位于 $\mathcal T$ 的内部，且 $i_0<K$。
理论策略从初始正长度阶段起即采用 $q=1$，到达安全集后取 $q=0$。
本例检验没有初始等待阶段的 $1\to0$ 结构，
而不是等待感染达到容量后才开始控制。

算例 E4 和 E5 均从容量线出发，但 $s_0$ 分别位于 $s_B$ 两侧：
\[
 0.70>s_B,\qquad 0.50<s_B.
\]
E4 属于容量等待弧 $\mathcal B$，理论上先采用 $q_B$ 保持容量，
在 $s=s_B$ 时转为完全跟踪；E5 则属于立即完全跟踪区，
理论上应直接采用 $q=1$ 并离开容量线。
因此，两者对应
\[
 \mathrm{E4}:\quad q_B\to1\to0,
 \qquad
 \mathrm{E5}:\quad 1\to0.
\]
这组对照说明，是否维持容量不能仅由 $i_0=K$ 决定，
还取决于状态在内生切换阈值两侧的位置。

另一个比较是 E2 与 E4 在容量退出时的状态。
两例到达该事件的时间不同，但理论上均经过同一状态 $(s_B,K)$；
从该点起的后续解析轨道相同。
因此可以在相平面中比较其数值退出位置，而不要求不同算例的时间序列重合。

\begin{figure}[H]
\centering
\NumFigureOrBlank{FigN3_boundary_tracking.pdf}{5.0cm}
\caption{立即完全跟踪与容量边界初值的对照。
从左到右为 E3、E4、E5，上排为 $q(t)$，下排为 $i(t)$。
解析参考与 OpenOCL 结果的比较重点为初始控制阶段、容量平台是否出现，以及解除位置。}
\label{fig:numerical-boundary-tracking}
\end{figure}

% EXPERIMENT_REQUIRED: observed initial stages and E2/E4 numerical exit states.
\NumFindingParagraph{\NumFindingTrackingBoundary}

\subsection{成本对照与结构总结}

表\ref{tab:numerical-scenarios}同时保留解析成本与直接配点成本，
用少量定量信息补充策略结构的比较。
其中解析成本依据各段显式成本求得，而数值成本依据式
\eqref{eq:numerical-cost-sum}计算。
比较不要求孤立切换时刻的控制代表值一致，
重点是各正长度阶段及其连接状态是否与理论预测相符。

% EXPERIMENT_REQUIRED: summarize supported structures only after all runs.
\NumFindingParagraph{\NumOverallFinding}
\FloatBarrier
